\documentclass{article}

\usepackage{amsmath,amssymb}
\usepackage{xurl}
\usepackage[hidelinks]{hyperref}

% Shared commands for notes in docs/paper-gaps/.

\DeclareUnicodeCharacter{2080}{\ifmmode{}_{0}\else\textsubscript{0}\fi}
\DeclareUnicodeCharacter{2081}{\ifmmode{}_{1}\else\textsubscript{1}\fi}
\DeclareUnicodeCharacter{2082}{\ifmmode{}_{2}\else\textsubscript{2}\fi}
\DeclareUnicodeCharacter{2099}{\ifmmode{}_{n}\else\textsubscript{n}\fi}

\newcommand{\Lean}{\textsc{Lean}}
\ExplSyntaxOn
\NewDocumentCommand{\leanid}{m}{
  \ifmmode
    \textsf{\detokenize{#1}}
  \else
    \group_begin:
    \urlstyle{sf}
    \tl_set:Nn \l_tmpa_tl {#1}
    \tl_replace_all:Nnn \l_tmpa_tl {\_} {_}
    \exp_args:NV \path \l_tmpa_tl
    \group_end:
  \fi
}
\ExplSyntaxOff
\providecommand{\tr}{\operatorname{tr}}
\newcommand{\tnleanIssueBase}{https://github.com/LionSR/TNLean/issues/}
\newcommand{\tnleanPrBase}{https://github.com/LionSR/TNLean/pull/}
\newcommand{\tnleanDocBase}{https://sirui-lu.com/TNLean/docs/}
\newcommand{\ghissue}[1]{\href{\tnleanIssueBase#1}{GitHub issue \##1}}
\newcommand{\ghpr}[1]{\href{\tnleanPrBase#1}{PR \##1}}
\newcommand{\doclink}[1]{\href{\tnleanDocBase#1.html}{\path{#1}}}

% Dirac notation and the identity operator, as in blueprint/src/macros/common.tex.
% \providecommand keeps note-local definitions authoritative.
\usepackage{dsfont}
\providecommand{\ket}[1]{|#1\rangle}
\providecommand{\bra}[1]{\langle #1|}
\providecommand{\braket}[2]{\langle #1 | #2 \rangle}
\providecommand{\ketbra}[2]{|#1\rangle\!\langle #2|}
\providecommand{\Id}{\mathds{1}}
\providecommand{\one}{\mathds{1}}
\providecommand{\C}{\mathbb{C}}
\providecommand{\MN}[1]{M_{#1}(\C)}
\providecommand{\MD}{M_D(\C)}

% Verdict marker: \gapnote{<kind>}{<status>}. Kind is the policy
% classification (clarification, local-correction, scope-restriction,
% unfaithful, false-source, open-gap); status is open, wip (elimination
% underway), resolved, or historical. Machine-read by the site index;
% prints a verdict line.
\newcommand{\gapnote}[2]{\par\noindent\textbf{Verdict:} #1 (#2).\par}


% Fallbacks keep this author document compilable when it is built outside its
% source directory and a different command.tex is found first.
\providecommand{\leanid}[1]{\textsf{\detokenize{#1}}}
\providecommand{\doclink}[1]{%
    \href{https://sirui-lu.com/TNLean/docs/#1.html}{\path{#1}}}
\providecommand{\gapnote}[2]{%
    \par\noindent\textbf{Verdict:} #1 (#2).\par}

\title{Exchange of Domain-Wall Endpoints on Glued Half-Chains}
\author{}
\date{2026-10-01}

\begin{document}
\maketitle
\gapnote{scope-restriction}{open}

Let $U$ be a matrix product unitary of order two exchanging the injective
ground states $\lvert\psi_A\rangle$ and $\lvert\psi_B\rangle$, let $e_{AB}$ and
$e_{BA}$ be domain walls on which $U$ acts with the phases $c_{AB}$ and
$c_{BA}$, and let $O^{[i,j]}$ be the string of $U$ on the sites $i,\dots,j$
whose two open virtual legs are closed by the endpoint tensors built from left
inverses of the ground-state tensors \cite{GarreRubioSchuch2024DomainWall}.
Ref.~\cite{GarreRubioSchuch2024DomainWall} cuts such a string in the middle
and obtains the half-chain operator $O^{[i]}_x$: the left endpoint at $i$
followed by the tensors of $U$, with the bond of the operator string fixed to
the value $x$ at a cut on the right. It proves, for $i>j$,
\begin{align}
    O^{[j]}_x\,O^{[i]}_y\,\lvert\psi_A\rangle
        &= c_{AB}\,O^{[i]}_x\,O^{[j]}_y\,\lvert\psi_A\rangle,
    \label{eq:gs24hs_open_exchange}
\end{align}
as an identity between half-chain objects with the two open legs $x$ and $y$,
and remarks that such an object must be glued with the other half of a string
to give a state.\footnote{Source traceability:
\path{Papers/2405.00439/MPU-DW.tex}, lines 1427--1659; the remark on gluing is
at line 1489.}

\section{The formal statement}

The formalization glues both sides of (\ref{eq:gs24hs_open_exchange}) with
right endpoints at two sites $r_1\neq r_2$ to the right of $i$, the leg $x$
being closed at $r_1$ and the leg $y$ at $r_2$. The relation becomes an
identity of periodic states,
\begin{align}
    O^{[j,r_1]}\,O^{[i,r_2]}\,\lvert\psi_A\rangle
        &= c_{AB}\,O^{[i,r_1]}\,O^{[j,r_2]}\,\lvert\psi_A\rangle,
    \label{eq:gs24hs_glued_exchange}
\end{align}
and it is proved for both orders of $r_1$ and $r_2$, whenever the regions
between the four endpoints are longer than the buffer of the local
action.\footnote{Formal declarations:
\leanid{MPOTensor.GroupFamily.BlockActionData.wallString_mul_wallString_mulVec_mpv_swap_left_far},
\leanid{MPOTensor.GroupFamily.BlockActionData.wallString_mul_wallString_mulVec_mpv_swap_left_near}.
The buffer is recorded in Ref.~\cite{gap:gs24_domain_wall_nondegenerate}.}
For $r_2<r_1$ the left side is $c_{AB}c_{BA}$ times the periodic state with the
domain walls $e_{AB},e_{BA},e_{AB},e_{BA}$ at $j,i,r_2,r_1$, and the right
side is $c_{BA}$ times the same state; for $r_1<r_2$ the two sides are $c_{AB}$
and $1$ times the state with these walls at $j,i,r_1,r_2$.

\section{Comparison}

Equation (\ref{eq:gs24hs_glued_exchange}) follows from
(\ref{eq:gs24hs_open_exchange}) by contracting both sides with the same right
halves, so it is a consequence of the source statement; the converse does not
follow, since only right halves of domain-wall strings are contracted. The
open-leg identity (\ref{eq:gs24hs_open_exchange}) uses in addition the
fixed-point form of the action tensors, which writes $m+1$ sites of $U$ acting
on $A$ as one acted site, the right action tensor, $m-1$ sites of $B$, the
left action tensor and one acted site. The formal action tensors satisfy
only the reduction of $m$ acted sites to $B^m$ between the two action tensors,
so the half-chain objects are not available with open legs.

\textbf{Verdict: open scope restriction.} Equation
(\ref{eq:gs24hs_glued_exchange}) is the formalized content of the
single-endpoint exchange. The open-leg identity
(\ref{eq:gs24hs_open_exchange}) needs the fixed-point decomposition of the
acted chain as an additional property of the action tensors.

\bibliographystyle{alpha}
\bibliography{references}

\end{document}
